\subsection{Ricerca ibrida e fusione dei risultati}

La ricerca per sola similarità vettoriale (\emph{dense retrieval}), tuttavia, ha un punto debole noto: tende a
performare peggio dei metodi lessicali classici (in particolare \emph{BM25} \cite{robertson2009bm25}) proprio sui termini rari o molto specifici: un codice di corso, una sigla normativa, un nome
proprio, che un embedding generico fatica a distinguere ma che un confronto esatto di parole individua senza
ambiguità. La letteratura piu recente raccomanda perciò l'adozione della \emph{ricerca ibrida} (\emph{hybrid search}): eseguire in parallelo una ricerca densa e una ricerca lessicale sparsa,
e fonderne i risultati in un'unica classifica. Una delle tecniche di fusione più diffuse a questo scopo,  è la \emph{Reciprocal Rank Fusion} (RRF), proposta da Cormack et al. \cite{cormack2009rrf}: per ogni documento si guarda non il punteggio ottenuto nelle due ricerche, ma solo la sua \emph{posizione} in ciascuna classifica e si sommano gli inversi di queste posizioni per ottenere un punteggio finale. Un documento che compare ai primi posti in entrambe le classifiche ottiene così un punteggio alto, uno che compare in alto in una sola delle due un punteggio intermedio. Il vantaggio pratico di questo meccanismo è che permette di fondere le due classifiche senza dover rendere confrontabili i punteggi originali delle due ricerche, cosa altrimenti problematica, perché il punteggio di una ricerca densa e quello di una ricerca lessicale sono calcolati in modo diverso e non sono direttamente comparabili tra loro.